
\documentclass{article}
\usepackage{colortbl}
\usepackage{makecell}
\usepackage{multirow}
\usepackage{supertabular}

\begin{document}

\newcounter{utterance}

\centering \large Interaction Transcript for game `bbh', experiment `formal\_fallacies', episode 56 with Qwen3.6{-}27B.
\vspace{24pt}

{ \footnotesize  \setcounter{utterance}{1}
\setlength{\tabcolsep}{0pt}
\begin{supertabular}{c@{$\;$}|p{.15\linewidth}@{}p{.15\linewidth}p{.15\linewidth}p{.15\linewidth}p{.15\linewidth}p{.15\linewidth}}
    \# & \multicolumn{2}{c}{Player} && \multicolumn{2}{c}{Game Master} \\
    \hline

    \theutterance \stepcounter{utterance}  
    & & & \multicolumn{4}{p{0.6\linewidth}}{
        \cellcolor[rgb]{0.9,0.9,0.9}{
            \makecell[{{p{\linewidth}}}]{
                \texttt{\tiny{[P1$\langle$GM]}}
                \texttt{Distinguish deductively valid arguments from formal fallacies.} \\
\\ 
\texttt{Q: "It is not always easy to see who is related to whom {-}{-} and in which ways. The following argument pertains to this question: To begin with, Lesley is a close friend of Fernando. Moreover, being a close friend of Fernando or a schoolmate of Lowell is sufficient for being a great{-}grandfather of Leroy. It follows that Lesley is a great{-}grandfather of Leroy."} \\
\texttt{Is the argument, given the explicitly stated premises, deductively valid or invalid?} \\
\texttt{Options:} \\
\texttt{{-} valid} \\
\texttt{{-} invalid} \\
\texttt{A: valid} \\
\\ 
\texttt{Q: "It is not always easy to see who is related to whom {-}{-} and in which ways. The following argument pertains to this question: Whoever is not a great{-}grandfather of Clyde is a stepbrother of Brian. Being an ancestor of Dana is sufficient for not being a great{-}grandfather of Clyde. We may conclude: Everyone who is an ancestor of Dana is a stepbrother of Brian, too."} \\
\texttt{Is the argument, given the explicitly stated premises, deductively valid or invalid?} \\
\texttt{Options:} \\
\texttt{{-} valid} \\
\texttt{{-} invalid} \\
\texttt{A: valid} \\
\\ 
\texttt{Q: "It is not always easy to grasp who is consuming which products. The following argument pertains to this question: Every infrequent user of Paul Mitchell shampoo is either a rare consumer of Nioxin shampoo or a loyal buyer of Caress soap, or both. No regular consumer of Lush soap is a rare consumer of Nioxin shampoo and, in the same time, a loyal buyer of Caress soap. It follows that whoever is an infrequent user of Paul Mitchell shampoo is not a regular consumer of Lush soap."} \\
\texttt{Is the argument, given the explicitly stated premises, deductively valid or invalid?} \\
\texttt{Options:} \\
\texttt{{-} valid} \\
\texttt{{-} invalid} \\
\texttt{A: invalid} \\
\\ 
\texttt{Q: "Here comes a perfectly valid argument: First, everyone who is a half{-}brother of Ronny or a classmate of Sal is not a son of Ronald. Second, it is false that Santo is a son of Ronald. So, necessarily, Santo is a half{-}brother of Ronny or a classmate of Sal."} \\
\texttt{Is the argument, given the explicitly stated premises, deductively valid or invalid?} \\
\texttt{Options:} \\
\texttt{{-} valid } \\
\texttt{{-} invalid} \\
\texttt{A:} \\
            }
        }
    }
     \\ \\

    \theutterance \stepcounter{utterance}  
    & \multicolumn{4}{p{0.6\linewidth}}{
        \cellcolor[rgb]{0.561,0.737,0.561}{
            \makecell[{{p{\linewidth}}}]{
                \texttt{\tiny{[P1$\rangle$GM]}}
                \texttt{To determine the validity of the argument, let's break it down into logical form.} \\
\\ 
\texttt{**Premises:**} \\
\texttt{1. Everyone who is a half{-}brother of Ronny (\$H\$) or a classmate of Sal (\$C\$) is not a son of Ronald (\$\textbackslash{}neg S\$).} \\
\texttt{   {-} Logical Form: \$(H \textbackslash{}lor C) \textbackslash{}rightarrow \textbackslash{}neg S\$} \\
\texttt{2. It is false that Santo is a son of Ronald.} \\
\texttt{   {-} Logical Form: \$\textbackslash{}neg S\$ (for Santo)} \\
\\ 
\texttt{**Conclusion:**} \\
\texttt{{-} Santo is a half{-}brother of Ronny or a classmate of Sal.} \\
\texttt{   {-} Logical Form: \$H \textbackslash{}lor C\$ (for Santo)} \\
\\ 
\texttt{**Analysis:**} \\
\texttt{The argument structure is:} \\
\texttt{1. \$P \textbackslash{}rightarrow Q\$} \\
\texttt{2. \$Q\$} \\
\texttt{3. Therefore, \$P\$} \\
\\ 
\texttt{Where:} \\
\texttt{{-} \$P\$ = Santo is a half{-}brother of Ronny or a classmate of Sal} \\
\texttt{{-} \$Q\$ = Santo is not a son of Ronald} \\
\\ 
\texttt{This is a classic formal fallacy known as **Affirming the Consequent**. Just because the consequent (\$\textbackslash{}neg S\$) is true, it does not logically follow that the antecedent (\$H \textbackslash{}lor C\$) must be true. There could be other reasons why Santo is not a son of Ronald (e.g., he is a stranger, a cousin, etc.). The premises do} \\
            }
        }
    }
    & & \\ \\

    \theutterance \stepcounter{utterance}  
    & & & \multicolumn{2}{p{0.3\linewidth}}{
        \cellcolor[rgb]{0.9,0.9,0.9}{
            \makecell[{{p{\linewidth}}}]{
                \texttt{\tiny{[GM$|$GM]}}
                \texttt{Target: invalid} \\
            }
        }
    }
    & & \\ \\

    \theutterance \stepcounter{utterance}  
    & & & \multicolumn{2}{p{0.3\linewidth}}{
        \cellcolor[rgb]{0.9,0.9,0.9}{
            \makecell[{{p{\linewidth}}}]{
                \texttt{\tiny{[GM$|$GM]}}
                \texttt{game\_result = LOSE} \\
            }
        }
    }
    & & \\ \\

\end{supertabular}
}

\end{document}
